
Our experiments validate whether zero-shot evaluation reveals architectural compatibility before parameter-efficient fine-tuning investment. We design experiments to answer specific research questions connecting to our central claims.

\subsubsection{Research Questions}

\textbf{RQ1: Does Mamba-130M demonstrate task-dependent zero-shot compatibility?} We hypothesize that causal state-space models succeed on generation-aligned classification tasks but fail on tasks requiring bidirectional reasoning. Success is defined as above-random baseline accuracy; failure is below-random accuracy indicating systematic architectural bias.

\textbf{RQ2: Is below-random performance statistically significant or sampling noise?} We test whether observed failures reflect genuine architectural constraints or insufficient sample sizes using binomial significance tests.

\textbf{RQ3: Does zero-shot performance correlate with task structure?} We predict that task requirements (symmetric comparison versus unidirectional encoding) determine success/failure patterns independent of task difficulty or domain.

\subsubsection{Tasks and Datasets}

We evaluate on three GLUE benchmark tasks representing diverse reasoning requirements:

\textbf{QQP (Quora Question Pairs):} Binary paraphrase detection over question pairs. Requires symmetric comparison---both questions must inform similarity judgment equally. Random baseline: 50\%. We sample 100 validation examples.

\textbf{MNLI (Multi-Genre Natural Language Inference):} Three-way classification (entailment, neutral, contradiction) over premise-hypothesis pairs. Requires bidirectional reasoning to compare premise and hypothesis symmetrically. Random baseline: 33.3\%. We evaluate on 100 matched validation examples.

\textbf{SST-2 (Stanford Sentiment Treebank):} Binary sentiment classification over single sentences. Aligns with causal language modeling---left-to-right encoding aggregates context sufficient for classification. Random baseline: 50\%. We evaluate on 100 validation examples.

\begin{table*}[htbp]
\centering
\resizebox{\dimexpr 2\columnwidth+\columnsep\relax}{!}{%
\begin{tabular}{lllll}
\toprule
Task & Type & Classes & Baseline & Architectural Requirement \\
\midrule
QQP & Paraphrase Detection & 2 & 50\% & Symmetric comparison \\
MNLI & Natural Language Inference & 3 & 33.3\% & Bidirectional reasoning \\
SST-2 & Sentiment Classification & 2 & 50\% & Unidirectional encoding \\
\bottomrule
\end{tabular}
}
\end{table*}

\subsubsection{Model Configuration}

We use the pretrained checkpoint \texttt{state-spaces/mamba-130m-hf} from HuggingFace. The model employs selective state-space mechanisms with 24 layers and approximately 130M parameters, pretrained on standard language modeling corpora using causal autoregressive objectives.

For each task, we attach a randomly initialized linear classification head mapping final hidden states to class logits. No pretraining or fine-tuning of the head is performed---true zero-shot evaluation.

Inference protocol uses deterministic evaluation with seed 42, greedy decoding, and no prompt engineering or few-shot examples. We measure raw zero-shot capability without task-specific optimization.

\subsubsection{Evaluation Metrics}

Primary metric is accuracy (percentage of correct predictions), compared against random baseline to determine architectural compatibility.

Statistical significance is assessed via binomial tests. For binary tasks (QQP, SST-2), random baseline is 50\%; for MNLI, 33.3\%. We use $\alpha$ = 0.05 significance threshold.

Interpretation thresholds:
\begin{itemize}
\item Accuracy $\geq$ random baseline: Model demonstrates task capability (architecture compatible)
\item Accuracy < random baseline: Systematic failure indicating architectural incompatibility
\end{itemize}

\subsubsection{Implementation Details}

All experiments run on a single NVIDIA GPU using HuggingFace Transformers library. For each task: (1) load pretrained Mamba-130M checkpoint, (2) initialize task-specific classification head, (3) perform forward passes on 100 validation examples without gradient updates, (4) compute accuracy and binomial significance.

Total evaluation time is under 5 minutes per task, demonstrating the practicality of zero-shot architectural compatibility testing as a PEFT prerequisite.

\subsubsection{Experimental Validity}

Sample size of 100 examples per task provides sufficient statistical power to detect below-random performance. For QQP, observing 38/100 correct (versus expected 50/100 under null hypothesis) yields p = 0.003, confirming robust significance.

We use the standard publicly available Mamba checkpoint rather than selecting from multiple variants, preventing cherry-picking bias. Results reflect base architectural capabilities, not checkpoint-specific artifacts.

Zero-shot evaluation removes confounds from hyperparameter optimization. Observed failures cannot be attributed to suboptimal learning rates or batch sizes---the architecture either supports the task or does not.

All experiments use deterministic seeds and publicly available datasets and checkpoints for reproducibility.

\subsubsection{Baseline Comparisons}

We establish architectural compatibility through comparison to random baselines rather than other models. This design isolates architectural capability: below-random accuracy definitively indicates incompatibility independent of model size, pretraining data, or training duration. Comparisons to GPT-2 or other baselines would confound architectural effects with checkpoint quality differences.
